\begin{table}[H]
  \centering
  \caption{Parte de la variación del nivel de pobreza y de la entrada en pobreza que
  concentra la comunidad, pobreza monetaria e IPM, 2010 a 2016.}
  \label{tab:robustez_icc}
  \footnotesize
  \setlength{\tabcolsep}{5pt}
  \begin{tabular}{l ccc ccc}
    \toprule
    & \multicolumn{3}{c}{\textbf{Monetaria}} & \multicolumn{3}{c}{\textbf{IPM}} \\
    \cmidrule(lr){2-4}\cmidrule(lr){5-7}
    & \textbf{Total} & \textbf{Urbano} & \textbf{Rural} & \textbf{Total} & \textbf{Urbano} & \textbf{Rural} \\
    \midrule
    Nivel de pobreza, 2010 & 26.7 & 21.9 & 15.2 & 20.9 & 14.9 & 13.6 \\
    Nivel de pobreza, 2013 & 22.2 & 24.5 & 10.1 & 12.6 & 11.5 & 9.0 \\
    Nivel de pobreza, 2016 & 17.6 & 19.5 & 13.8 & 8.8 & 7.2 & 8.0 \\
    Entrada en pobreza, 2010$\to$2013 & 12.5 & 11.8 & 6.0 & 5.9 & 6.7 & 3.6 \\
    Entrada en pobreza, 2013$\to$2016 & 11.2 & 13.9 & 7.2 & 1.8 & 1.7 & 3.1 \\
    \bottomrule
  \end{tabular}
  \begin{minipage}{0.95\textwidth}
    \vspace{4pt}
    \footnotesize \textit{Nota:} Porcentaje de la variación que corresponde a la
    comunidad: correlación intraclase de un análisis de varianza de un factor, corregida
    por el tamaño desigual de las comunidades (Anexo~\ref{apx:variacion_dmsp}). Nivel de
    pobreza: todos los hogares de cada ola. Entrada en pobreza: hogares no pobres al
    comienzo de cada ventana, con o sin dato de DMSP-OLS. Se excluyen los hogares sin
    comunidad identificada.
    Fuente: cálculos propios.
  \end{minipage}
\end{table}
